\documentclass[11pt,a4paper]{article}

\usepackage[utf8]{inputenc}
\usepackage[T1]{fontenc}
\usepackage[french]{babel}
\usepackage{lmodern}
\usepackage{microtype}
\usepackage[margin=2.5cm,headheight=15pt]{geometry}
\usepackage{xcolor}
\usepackage{graphicx}
\usepackage{eso-pic} 
\usepackage{titlesec}
\usepackage{fancyhdr}
\usepackage{booktabs,tabularx,array}
\usepackage{listings}
\usepackage{hyperref}

\usepackage{tikz}
\usetikzlibrary{positioning,arrows.meta,shapes.geometric}
\tikzset{
  entite/.style={
    draw=accent, rounded corners=3pt, fill=accent!7,
    minimum width=3cm, minimum height=0.9cm,
    align=center, font=\sffamily\bfseries
  },
  association/.style={
    draw=discret, diamond, aspect=2, fill=white,
    minimum width=2cm, minimum height=0.8cm,
    align=center, font=\sffamily
  },
  lien/.style={draw=discret, -{Latex[length=2mm]}}
}

\newcommand{\TitreProjet}{Test du projet}
\newcommand{\SousTitre}{Conception et réalisation d'une base de données}
\newcommand{\Universite}{LSIN513 - Bases de données}
\newcommand{\Formation}{L3 Informatique}
\newcommand{\Annee}{2026--2027}
\newcommand{\AuteurUn}{Prénom Nom 1}
\newcommand{\AuteurDeux}{Prénom Nom 2}
\newcommand{\AuteurTrois}{Prénom Nom 3}

\newcommand{\FichierLogo}{UVSQ.jpg}



\newif\ifrapportfinal
\rapportfinaltrue








\definecolor{accent}{HTML}{175673}
\definecolor{encre}{HTML}{24343D}
\definecolor{discret}{HTML}{65727A}
\definecolor{fondcode}{HTML}{F3F6F7}
\definecolor{trait}{HTML}{D8E2E5}

\hypersetup{
  colorlinks=true,linkcolor=accent,urlcolor=accent,citecolor=accent,
  pdftitle={\TitreProjet},
  pdfauthor={\AuteurUn, \AuteurDeux, \AuteurTrois},
  pdfsubject={Rapport de projet de bases de données}
}
\urlstyle{same}
\setlength{\parindent}{0pt}
\setlength{\parskip}{0.55em}
\setlength{\emergencystretch}{2em}
\renewcommand{\arraystretch}{1.25}
\newcolumntype{Y}{>{\raggedright\arraybackslash}X}
\setcounter{tocdepth}{2}

\titleformat{\section}
  {\Large\sffamily\bfseries\color{encre}}
  {\color{accent}\thesection}{0.7em}{}
  [\vspace{0.2em}{\color{trait}\titlerule}]
\titleformat{\subsection}
  {\large\sffamily\bfseries\color{accent}}
  {\thesubsection}{0.7em}{}
\titlespacing*{\section}{0pt}{1.6em}{0.7em}
\titlespacing*{\subsection}{0pt}{1em}{0.3em}

\pagestyle{fancy}
\fancyhf{}
\fancyhead[R]{\small\sffamily\color{discret}\Formation}
\fancyfoot[L]{\small\sffamily\color{discret}Rapport de projet}
\fancyfoot[R]{\small\sffamily\color{accent}\thepage}
\renewcommand{\headrulewidth}{0.3pt}


\newcommand{\aCompleter}[1]{{\color{discret}\itshape #1}}



\newcommand{\emplacementFigure}[1]{
  \fcolorbox{trait}{fondcode}{
    \parbox[c][3cm][c]{0.93\linewidth}{
      \centering\sffamily\color{discret}#1}}}


\lstdefinestyle{sqlprojet}{
  language=SQL,
  morekeywords={DECLARE,BEGIN,END,EXCEPTION,RETURN,LOOP,ELSIF,
    DBMS_OUTPUT,PUT_LINE,RAISE_APPLICATION_ERROR},
  basicstyle=\small\ttfamily,
  keywordstyle=\bfseries\color{accent},
  commentstyle=\itshape\color{discret},
  stringstyle=\color{brown},
  backgroundcolor=\color{fondcode},
  frame=single,rulecolor=\color{trait},framerule=0.4pt,
  numbers=left,numberstyle=\scriptsize\color{discret},numbersep=9pt,
  xleftmargin=1.7em,framexleftmargin=0.4em,
  breaklines=true,columns=fullflexible,keepspaces=true,
  showstringspaces=false,tabsize=2,captionpos=b,
  aboveskip=1em,belowskip=0.7em,
  literate={é}{{\'e}}1 {è}{{\`e}}1 {ê}{{\^e}}1 {ë}{{\"e}}1
    {à}{{\`a}}1 {â}{{\^a}}1 {î}{{\^i}}1 {ï}{{\"i}}1
    {ô}{{\^o}}1 {ù}{{\`u}}1 {û}{{\^u}}1 {ü}{{\"u}}1
    {ç}{{\c c}}1 {É}{{\'E}}1 {À}{{\`A}}1 {œ}{{\oe}}1
}
\lstset{style=sqlprojet}
\renewcommand{\lstlistingname}{Extrait}

\begin{document}


\begin{titlepage}
  
  
  
  \IfFileExists{\FichierLogo}{
    \AddToShipoutPictureFG*{
      \AtPageUpperLeft{
        \put(\LenToUnit{\dimexpr\paperwidth-2.5cm\relax},\LenToUnit{-2.2cm}){
          \makebox[0pt][r]{
            \raisebox{-\height}[0pt][0pt]{
              \includegraphics[width=3cm,height=3cm,keepaspectratio]{\FichierLogo}}}}}
    }
  }{}
  \sffamily
  {\large\color{encre}\Universite\par}
  {\color{discret}\Formation\enspace - \enspace\Annee\par}
  \vspace{0.7cm}
  {\color{accent}\rule{\linewidth}{1.5pt}\par}

  \vspace{2.3cm}
  {\small\bfseries\color{accent}RAPPORT DE PROJET\par}
  \vspace{0.4cm}
  {\fontsize{30}{36}\selectfont\bfseries\color{encre}\TitreProjet\par}
  \vspace{0.55cm}
  {\Large\color{discret}\SousTitre\par}

  \vfill
  \begin{minipage}[t]{0.54\textwidth}
    {\small\bfseries\color{accent}ÉQUIPE\par}
    \vspace{0.2cm}
    \AuteurUn\par
    \AuteurDeux\par
    \AuteurTrois\par
    \vspace{0.15cm}
  \end{minipage}\hfill
  \vspace{0.8cm}
\end{titlepage}

\pagenumbering{roman}
\tableofcontents
\clearpage
\pagenumbering{arabic}




\section{Phase 1 : analyse et modélisation}




\subsection{Cahier des charges}

\aCompleter{Décrire le cahier des charges, c’est-à-dire la description des besoins de votre application
de bases de données.}

% \subsubsection{Contexte et objectif}
% \aCompleter{Présenter le sujet retenu, le besoin auquel il répond et les
% utilisateurs concernés. Résumer l'objectif du projet en quelques phrases.}

% \subsubsection{Périmètre et fonctionnalités}
% \aCompleter{Lister les fonctionnalités prévues et préciser les limites
% du projet. Adapter le tableau à votre sujet.}

% \begin{tabularx}{\linewidth}{@{}lYY@{}}
%   \toprule
%   \textbf{Référence} & \textbf{Fonctionnalité} & \textbf{Résultat attendu} \\
%   \midrule
%   F01 & \aCompleter{À compléter} & \aCompleter{À compléter} \\
%   F02 & \aCompleter{À compléter} & \aCompleter{À compléter} \\
%   \bottomrule
% \end{tabularx}

\subsection{Modélisation}
\aCompleter{Donner le schéma Entité/Association détaillé modélisant les besoins définis
dans le cahier des charges. On rappelle qu’un schéma E/A est indépendant d’un modèle de
données.}

\begin{figure}[htbp]
  \centering
  \emplacementFigure{Emplacement du schéma conceptuel}
  \caption{Modèle conceptuel de la base de données.}
  \label{fig:modele}
\end{figure}

\subsection{Contraintes d'intégrité}
\aCompleter{Définition les contraintes d'intégrité qui ne peuvent pas être
exprimées sur le modèle conceptuel. Mettre l’accent sur la rédaction en langage naturel de
contraintes d’intégrité complexes.}

\subsection{Droits, vues et confidentialité}
\aCompleter{Décrire les différents utilisateurs ou groupes
d’utilisateurs de votre base de données, ainsi que les droits qu’ils devraient avoir en vous
référant aux concepts décrits au-dessus.}

\begin{tabularx}{\linewidth}{@{}YYY@{}}
  \toprule
  \textbf{Utilisateur / groupe} & \textbf{Informations accessibles} & \textbf{Actions autorisées} \\
  \midrule
  \aCompleter{À compléter} & \aCompleter{À compléter} & \aCompleter{À compléter} \\
  \bottomrule
\end{tabularx}

\subsection{Requêtes}
\aCompleter{Qu’est-ce qui intéresse les utilisateurs de votre base de données? Définir en
langage naturel une vingtaine de questions relativement complexes en langage naturel
nécessitant des calculs d'agrégats, sélection, projection, jointure, requêtes imbriquées,
intersection, union, différence et division, jointure externe etc..}



\begin{description}
  \item[Q01.] \aCompleter{Première question à définir.}
  \item[Q02.] \aCompleter{Deuxième question à définir.}
  
\end{description}

\subsection{Traduction du schéma E/A en relationnel}
\aCompleter{Traduction du schéma E/A en relationnel. Lister le nom des tables avec le nom de leurs
colonnes, en indiquant en gras les clés primaires et les clés étrangères soulignées.}


\begin{tikzpicture}[node distance=2.3cm]
  \node[entite] (truc) {Truc};
  \node[association, right=of truc] (compat) {est compatible};
  \node[entite, right=of compat] (machin) {Machin};

  \draw[lien] (truc) -- node[above, font=\small]{1,1} (compat);
  \draw[lien] (compat) -- node[above, font=\small]{0,n} (machin);
\end{tikzpicture}

\end{document}